\section{The native grading of the original multicone algebra}

\begin{lemma}[The original algebra has its literal total homogeneous grading]
\label{lem:fc-native-model-total-grading}
\leanok
\lean{EnvelopeProperty.MultiModel.homogeneousComponent_mem_model, EnvelopeProperty.MultiModel.nativeTotalGrading, EnvelopeProperty.MultiModel.mem_nativeTotalGrading, EnvelopeProperty.MultiModel.nativeTotalGrading_gradedMonoid, EnvelopeProperty.MultiModel.nativeTotalGrading_isInternal, EnvelopeProperty.MultiModel.nativeTotalGradedAlgebra}
\uses{def:gl-model}
For every original integral multiprojective model, its SAME original
subalgebra of $K[T_1,\ldots,T_r]$ has a native total-degree graded-algebra
structure. Its degree-$d$ submodule consists precisely of the original
polynomials homogeneous of total degree $d$.
\end{lemma}
\begin{proof}
Every total homogeneous component is the finite sum of its original
multihomogeneous monomials. The defining gradedness of the original model
puts each such monomial, and therefore their sum, in the original algebra.
Multiplication and the unit preserve the stated degrees in the ambient
polynomial algebra. The independence of homogeneous pieces is inherited
from the actual ambient direct-sum decomposition by the injective original
subalgebra inclusion. Every original polynomial is the finite sum of its
actual total homogeneous components, so the pieces span the whole
original algebra. This internal direct sum gives its native graded-algebra
structure.
\end{proof}

\begin{lemma}[The native inclusion and components retain the original polynomial]
\label{lem:fc-native-model-total-components}
\leanok
\lean{EnvelopeProperty.MultiModel.nativeTotalInclusion, EnvelopeProperty.MultiModel.nativeTotalGrading_decompose_val}
\uses{lem:fc-native-model-total-grading}
The original subalgebra inclusion is a native graded algebra morphism,
and every degree component in the new native decomposition is the SAME
original ambient homogeneous component.
\end{lemma}
\begin{proof}
Membership in each piece is precisely ambient homogeneity, so the
original inclusion preserves every degree. Naturality of native direct-sum
decomposition under this graded inclusion identifies its degree-$d$
component with the actual polynomial homogeneous component.
\end{proof}

For rank one this grading supplies the intended ordinary projective
scheme. For higher rank, total-degree Proj is not identified with the
original multiprojective model: the intended route first takes the
original positive diagonal Veronese and then its rank-one grading.
